\begin{table*}[t]
\centering
\footnotesize
\begin{tabular}{l|p{.9\textwidth}}
\textbf{RL} & \textbf{SFT Output} \\\hline

1 & . These have to brain with for to. The, the brain, and the. This and to. These and, and the cells. The to a and to the heart that cells to is cells, are the their their. In, it the heart of the brain this. The. Next of are cells the brain the heart the and by. But the that the cells to to this the et al to the cells. These, and to the environment, a. We the mechanism. The, and in,. In the body and are is. This cells, the body, and of the heart to the body to of a a of, the brain. The, the. \\\hline

3 & brain a a of to the to, they. However. \\\hline

5 & , the that the cells to the of the the of of the visual in and of the brain are, the the body. The. In, but the time to the eyes the brain and the brain cells. The their are to the brain. Here, and the cells. \\\hline

7 & enerative their body to, the two the, in the body of, the their into. We from the cell of, the into the by their cells. \\\hline

9 & brain. In. \\\hline

\end{tabular}
\caption{Outputs generated by SFT model and the intended Readability Level (RL)}
\label{tab:sft-examples}
\end{table*}


\begin{table*}[t]
\centering
\begin{tabular}{l|p{.9\textwidth}}
\textbf{Readability Level} & \textbf{LoRA Output} \\\hline

1 & Small open reading frames (smORFs) are small stretches of DNA that code for proteins, but their translation into proteins has been difficult to detect. This is because they are too short to produce the high conservation scores that are used to predict which stretches of DNA code for proteins. To overcome this problem, Kondziella et al. have developed a new method called Poly-Ribo-Seq that uses a combination of biochemistry and sequencing to detect smORFs that are being translated into proteins. Using this method, they were able to detect smORFs in Drosophila cells that were not detected by other methods. These smORFs were found to be translated at similar levels to canonical proteins and had the same characteristics as canonical proteins. These findings suggest that many smORFs in Drosophila are being translated into proteins and may have important functions in the cell. Future experiments will be needed to confirm this and to determine the functions of these proteins. \\ \hline

3 & The discoidin domain receptors (DDRs) are a family of receptor tyrosine kinases (RTKs) that bind to collagen and play a key role in tissue development, cell migration, and disease progression. DDRs are composed of an extracellular region that binds to collagen, a transmembrane domain and a cytoplasmic region that contains the catalytic kinase domain. Unlike other RTKs, DDRs are constitutively dimerised in the absence of ligand and activate downstream signalling through a mechanism that is not well understood. Here, we show that DDR1 activation by collagen involves phosphorylation of the juxtamembrane and activation loop tyrosines in trans between dimers. Phosphorylation between DDR1 dimers requires specific contacts between the transmembrane domains but not between the extracellular domains. These findings define a unique activation mechanism for DDR1. \\ \hline

5 & Stem cell biology is a rapidly evolving field that is critical for our understanding of how tissues and organs form during development and how they are maintained in the adult. In the fruit fly Drosophila melanogaster, the lymph gland is a multi-lobed organ that produces blood cells throughout the larval stages. This organ contains a stem cell niche that maintains a pool of undifferentiated blood cell progenitors. While the adult hematopoietic hub is known to be the final site of hematopoiesis, the mechanism of blood cell formation in larvae has been largely unexplored. Here we identify a population of cells in the early larval lymph gland that express the Notch signaling pathway. These cells are multipotent, dividing and give rise to the domeless expressing blood cell progenitors that populate the lymph gland. Notch signaling is required for their maintenance and relies on Dpp signaling from the niche. These results reveal the presence of hematopoietic stem cells in the early larval lymph gland and establish Drosophila as a model for studying hematopoietic stem cell biology. The findings also shed light on the evolutionary conservation of hematopoietic stem cell biology between Drosophila and vertebrates. \\ \hline

7 & The malaria parasite Plasmodium causes a disease that kills more than a million people every year. To infect the human body, the parasite must first invade the liver cells. How this invasion takes place is still not well understood. Silvie et al. have now investigated how Plasmodium parasites enter human liver cells. The experiments show that two host proteins, CD81 and SR-BI, can be used by the parasite to invade the liver cells. The parasite uses different proteins to enter the liver cells of different species. The experiments also show that a parasite protein called P36 can bind to the host proteins and trigger the invasion of the liver cells. The experiments provide new insights into the invasion process and could ultimately help to develop new treatments for malaria. \\ \hline

9 & A large family of proteins called TRIM5 is thought to be a key part of the body's natural defense against viruses. These proteins prevent certain viruses from multiplying by binding to the proteins that form the outer shell of the virus. To find out how TRIM5 proteins work, researchers have made a detailed 3D image of a TRIM5 protein bound to a virus. The image shows that the protein forms a hexagonal pattern around the virus, with each hexagon made up of 12 smaller protein molecules. The next step will be to find out how the proteins interact with the virus to prevent it from multiplying. These findings could also help to develop new treatments for viruses. \\ \hline

\end{tabular}
\caption{Outputs generated by LoRA model}
\label{tab:lora-examples}
\end{table*}

\begin{table*}[t]
\centering
\begin{tabular}{l|p{.9\textwidth}}
\textbf{Readability Level} & \textbf{DSPy Output} \\\hline

1 & ""\\\hline

3 & ""\\\hline

5 & ""\\\hline

7 & ""\\\hline

9 & ""\\ \hline

\end{tabular}
\caption{Outputs generated by DSPy (empty because of length failure)}
\label{tab:dspy-examples}
\end{table*}